% !TEX root = ../aa_SoM_Chapters.tex

% ö = \%c3\%b6
% ä = \%C3\%A4

\xdef\sectiontitle{\IIIsectiontitle} %aiheuttama

%%%%%%%%%%%%%%%
\begin{frame}[label=3_6_a]%[label=3_6_TOC1]
\frametitle{\texorpdfstring{\culine[\sectiontitleulinecolor]{\sectiontitle}}{\sectiontitle} }

\vspace{1cm}

\begin{itemize}[leftmargin=0.5cm,itemsep=3mm]
    \item[3.1] \hyperlink{3_1_a}{\IIIaSubsectiontitle}
    \item[3.2] \hyperlink{3_2_a}{\IIIbSubsectiontitle}	
    \item[3.3] \hyperlink{3_3_a}{\IIIcSubsectiontitle}	
    \item[3.4] \hyperlink{3_4_a}{\IIIdSubsectiontitle}
    \item[3.5] \hyperlink{3_5_a}{\IIIeSubsectiontitle}
    \item[3.6] \hyperlink{3_6_a}{\CDAlert[red]{\IIIfSubsectiontitle}}
    \item[3.7] \hyperlink{3_7_a}{\IIIgSubsectiontitle}
\end{itemize}

%\endofslide 
\end{frame}
%%%%%%%%%%%%%%%

\xdef\sectiontitle{\IIIfSubsectiontitle} 


%%%%%%%%%%%%%%%
\begin{frame}
\frametitle{\texorpdfstring{\culine[\sectiontitleulinecolor]{\sectiontitle}}{\sectiontitle} }
\framesubtitle{Spontaneous and induced nuclear reactions}
%On Spontaneous and Induced Nuclear Reactions

\begin{itemize}
	\item Nuclear reactions can be:
\begin{itemize}
	\item \CDAlert[fuchsia]{Spontaneous}: \appear{such as in the case of radioactivity, as described above.} \appear{Then the $Q$-value of the reaction needs to be positive for the reaction to happen}
\item \CDAlert[green]{Induced}: \appear{where we have an incoming (high-energy) particle} \appear{to collide on the target nucleus and to enable the reaction}

\end{itemize}

\item When the target nucleus is bombarded with a projectile particle\appear{, the collision may be:}
\begin{itemize}
	\item \CDAlert[red]{Elastic}: Elastic scattering is reflection from the edge of the potential well of target nucleus \appear{$\oldrightarrow$ the incoming particle practically \Alert[red]{changes momentum} and direction}
\item \CDAlert[blue]{Inelastic}: where a \Alert[blue]{direct interaction} occurs\appear{, and as a result, \Alert[blue]{usually an intermediate 'compound nucleus' is formed first}}\appear{, before finally decaying into reaction products}

\end{itemize}

\end{itemize}

%\xdef\loc{south east}
%\begin{tikzpicture}[remember picture,overlay] % anchor=south
%    \node (A) [yshift=-0.25*\figYshift,xshift=\figXshift, anchor=\loc] at (current page.\loc) {\includegraphics[width=7cm]{Figures_booklet/SOM_12_Nuclear_physics_figures/SOM_Nuclear_physics_Figure_1.png}}; % 8cm max height
%\end{tikzpicture}

\endofslide 
%\endofsection
\end{frame}
%%%%%%%%%%%%%%%

%%%%%%%%%%%%%%%
\begin{frame}
\frametitle{\texorpdfstring{\culine[\sectiontitleulinecolor]{\sectiontitle}}{\sectiontitle} }
\framesubtitle{Inelastic scattering \& Induced nuclear reaction}

%  
\begin{itemize}
	\item Inelastic scattering often occurs at intermediate energy levels\appear{, still less than ${\sim}140 \mathrm{\;MeV}$ }
	\implies{ \CDAlert[blue]{no extra particles} \Alert[blue]{will yet form} (such as mesons)}
	
	 \item The collision of the \CDAlert[blue]{target nucleus} \appear{($X$)} \appear{with a \CDAlert[tuniblue]{projectile}} \appear{($x$)} \appear{creates a compound nucleus}
	 
	 \item  The compound nucleus decays into \CDAlert[red]{product nucleus} \appear{($Y$)}\appear{, often accompanied with an \CDAlert[tunired]{ejected particle} ($y$) or particles}
	 
	 \item The \CDAlert[fuchsia]{general reaction scheme} is either written as
\[
 x ~+~ X \quad \rightarrow \quad \appear{Y ~+~ y ~+~ Q}  
\]
\appear{or slightly more concisely using \Alert[fuchsia]{parentheses convention}:}
\[
 \Alert[blue]{target nucleus} \appear{\;(\; } \appear{\Alert[tuniblue]{projectile}} \appear{\;,\;} \appear{\Alert[tunired]{ejected particle}} \appear{\;)\;} \appear{\Alert[red]{product nucleus}}
\]
\appear{i.e. as} \vspace{-4mm}
\[
 X(x, y) Y
 \]
\end{itemize}

%\xdef\loc{south east}
%\begin{tikzpicture}[remember picture,overlay] % anchor=south
%    \node (A) [yshift=-0.25*\figYshift,xshift=\figXshift, anchor=\loc] at (current page.\loc) {\includegraphics[width=7cm]{Figures_booklet/SOM_12_Nuclear_physics_figures/SOM_Nuclear_physics_Figure_1.png}}; % 8cm max height
%\end{tikzpicture}

\endofslide 
%\endofsection
\end{frame}
%%%%%%%%%%%%%%%

%%%%%%%%%%%%%%%
\begin{frame}
\frametitle{\texorpdfstring{\culine[\sectiontitleulinecolor]{\sectiontitle}}{\sectiontitle} }
%\framesubtitle{Inelastic scattering \&  Induced nuclear reaction}
\framesubtitle{Collisions and Q-value of the reaction}

\begin{itemize}
	\item Whatever the nuclear reaction will be\appear{, it needs to conserve energy.} \appear{For the nuclear reaction $X(x, y) Y$}\appear{, the energy released in the reaction is called the \CDAlert[fuchsia]{$Q$-value}:}
\[
x ~+~ X \quad \rightarrow \quad \appear{Y ~+~ y ~+~ Q}  \appear{\;, \text{ where }} \qquad \appear{Q=}\appear{(m_{x}+m_{X}-m_{y}-m_{Y})} \appear{c^{2}}
\]
\item $Q>0$: \appear{\CDAlert[red]{exothermic} reaction} \appear{(initial particles are heavier than products) }
\item $Q<0$: \appear{\CDAlert[blue]{endothermic} reaction} \appear{(initial particles are lighter)}

\item Examples: \vspace{-6mm}
$$
\begin{aligned}
 \appear{\mathrm{n} ~+~ { }^{1} \mathrm{H} \quad &\oldrightarrow \quad { }^{2} \mathrm{H} ~+~ \gamma ~+~ 2.22 \mathrm{\;MeV}} \quad  \appear{\text{(Exothermic)}} \\
\appear{\gamma ~+~ { }^{2} \mathrm{H} \quad &\oldrightarrow \quad { }^{1} \mathrm{H} ~+~ \mathrm{n} ~-~ 2.22 \mathrm{\;MeV}} \quad \appear{\text{(Endothermic)}}
\end{aligned}
$$

\item Endothermic reaction requires a \CDAlert[blue]{threshold energy} to be supplied %to the system 

\item According to the above\appear{, the threshold energy is $|Q|$.} \appear{But this is valid only for the \CDAlert[fuchsia]{center-of-mass (CM) coordinate frame}}
\end{itemize}

%\xdef\loc{south east}
%\begin{tikzpicture}[remember picture,overlay] % anchor=south
%    \node (A) [yshift=-0.25*\figYshift,xshift=\figXshift, anchor=\loc] at (current page.\loc) {\includegraphics[width=7cm]{Figures_booklet/SOM_12_Nuclear_physics_figures/SOM_Nuclear_physics_Figure_1.png}}; % 8cm max height
%\end{tikzpicture}

\endofslide 
%\endofsection
\end{frame}
%%%%%%%%%%%%%%%

%%%%%%%%%%%%%%%
\begin{frame}
\frametitle{\texorpdfstring{\culine[\sectiontitleulinecolor]{\sectiontitle}}{\sectiontitle} }
\framesubtitle{Kinetics}


\seminormal
% 6.5 cm = midway, 10 cm = 3/4 
%\vspace{2.5mm}
\begin{minipage}{9.2cm}
\begin{itemize}[itemsep=2mm]
	\item We examine first the reaction in the \CDAlert[blue]{CM frame} \appear{with non-relativistic case} \appear{of $x$ of mass $m$ (velocity $v_p$) incident on $X$ of mass $M$ (velocity $v_c$):}
	\appear{\xdef\loc{north east}%
\begin{tikzpicture}[remember picture,overlay] % anchor=south
    \node (A) [yshift=0.75*\figYshift,xshift=\figXshift, anchor=\loc] at (current page.\loc) {\includegraphics[width=3.8cm]{Figures_booklet/SOM_Tipler_Nuclear_physics_figures/SOM_Tipler_Nuclear_physics_figure_50_cropped_a}};% 8cm max height
\end{tikzpicture}}
$$
\begin{aligned}
\appear{\vec{p}_{m}+\vec{p}_{M}}\appear{=0} \qquad &\Rightarrow \appear{\qquad |\vec{p}_{m} |= |\vec{p}_{M} |} \equiv \appear{p} \\
  \qquad &\Rightarrow \qquad \appear{m v_p=M V_c} \qquad \Rightarrow \qquad \appear{\Frac{m}{M}=\Frac{V_c}{v_p} } \\
  \appear{E_{CM}}  \equal \frac{p^{2}}{2 m}&  \plus \frac{p^{2}}{2 M} \equal \frac{p^{2}}{2}  \LP \frac{m+M}{m M} \,\RP  \appear{\,,} \qquad \appear{\text{CM frame}}
\end{aligned}
$$
\end{itemize}
\end{minipage}
\vspace{2.0mm}

\begin{itemize}[itemsep=2mm]

\item While on the \CDAlert[red]{laboratory frame},%
\appear{\xdef\loc{south east}%
\begin{tikzpicture}[remember picture,overlay] % anchor=south
    \node (A) [yshift=-0.15*\figYshift,xshift=\figXshift, anchor=\loc] at (current page.\loc) {\includegraphics[width=3.8cm]{Figures_booklet/SOM_Tipler_Nuclear_physics_figures/SOM_Tipler_Nuclear_physics_figure_50_cropped_b}};% 8cm max height
\end{tikzpicture}}
 \appear{where the velocity of the incident particle $\mathrm{x}$ is $v_p+V_c$}\appear{, towards the target $M$}\appear{, which is at rest}\appear{, we obtain the \CDAlert[red]{momentum}:}
\[
p_{lab}  \equal \appear{m(v_p+V_c)}  \equal \appear{m v_p} \LP \appear{1+}\frac{V_c}{v_p} \,\RP  \equal \appear{m v_p} \LP \appear{1+\Frac{m}{M}} \,\RP  \equal  \appear{p} \LP \appear{1+\Frac{m}{M}} \,\RP \hspace{3.5cm}
\] 
\appear{and for \CDAlert[red]{energy}}\vspace{-6.5mm}
\[
E_{lab}  \equal \frac{p_{lab}^{2}}{2 m}  \equal \frac{p^{2}}{2 m} \LP \frac{m+M}{M} \,\appear{\RP^{2}}  \equal \frac{m+M}{M} \appear{E_{C M}} \hspace{3.5cm}
\]
\end{itemize}



\endofslide 
%\endofsection
\end{frame}
%%%%%%%%%%%%%%%

%%%%%%%%%%%%%%%%
%\begin{frame}
%\frametitle{\texorpdfstring{\culine[\sectiontitleulinecolor]{\sectiontitle}}{\sectiontitle} }
%\framesubtitle{Kinetics}
%
%\seminormal
%\begin{itemize}
%	\item Thus, we get the threshold energy for the endothermic reaction in lab frame:
%\[
%E_{\text{threshold}}  \equal  \frac{m+M}{M} \appear{|Q|}
%\]
%\end{itemize}
%
%%\xdef\loc{south east}
%%\begin{tikzpicture}[remember picture,overlay] % anchor=south
%%    \node (A) [yshift=-0.25*\figYshift,xshift=\figXshift, anchor=\loc] at (current page.\loc) {\includegraphics[width=7cm]{Figures_booklet/SOM_12_Nuclear_physics_figures/SOM_Nuclear_physics_Figure_1.png}}; % 8cm max height
%%\end{tikzpicture}
%
%\endofslide 
%%\endofsection
%\end{frame}
%%%%%%%%%%%%%%%%

%%%%%%%%%%%%%%%
\begin{frame}
\frametitle{\texorpdfstring{\culine[\sectiontitleulinecolor]{\sectiontitle}}{\sectiontitle} }
\framesubtitle{Threshold energy and Example}

\seminormal
\begin{itemize}
	\item Thus, we get the threshold energy for the endothermic reaction in lab frame:
\[
E_{\text{threshold}}  \equal  \frac{m+M}{M} \appear{|Q|}
\]

	\item As an \CDAlert[fuchsia]{example}\appear{, let us calculate the minimum kinetic energy of protons incident on \CDAlert[fuchsia]{${ }^{13} C$ nuclei at rest in the laboratory}}\appear{, that will produce the endothermic reaction:} 
	\[
	\quad { }_{\;6}^{13} C  ~\plus~  \appear{p} \quad \rightarrow \quad \appear{{ }_{\;7}^{13}N} \appear{ ~+~ n}  \qquad \quad \appear{i.e.} \qquad \quad \appear{{ }_{\;6}^{13} C}\appear{(p} \appear{,} \appear{n)} \appear{{ }_{\;7}^{13}N}
	\]

\item The $Q$ value for the reaction is:
$$
\begin{aligned}
\appear{Q&=} \LP \appear{m_{\text {initial }}-m_{\text{final}}} \,\RP  \appear{c^{2}} \\
&\appear{=(13.003355 \;u+1.007825 \;u-13.005738 \;u-1.008665 \;u)} \appear{\times 931.5 \mathrm{\;MeV} / u} \\
&\appear{=-3.00 \mathrm{\;MeV}} \\
\appear{E_{\text{threshold}}&=}\frac{m+M}{M} \appear{|Q|}  \equal  \frac{1.007825 \;u+13.003355 \;u}{13.003355 \;u} \appear{\times} \appear{3.00 \mathrm{\;MeV}}  \equal \appear{3.23 \mathrm{\;MeV}}
\end{aligned}
$$
\end{itemize}

%\xdef\loc{south east}
%\begin{tikzpicture}[remember picture,overlay] % anchor=south
%    \node (A) [yshift=-0.25*\figYshift,xshift=\figXshift, anchor=\loc] at (current page.\loc) {\includegraphics[width=7cm]{Figures_booklet/SOM_12_Nuclear_physics_figures/SOM_Nuclear_physics_Figure_1.png}}; % 8cm max height
%\end{tikzpicture}

\endofslide 
%\endofsection
\end{frame}
%%%%%%%%%%%%%%%

%%%%%%%%%%%%%%%
\begin{frame}
\frametitle{\texorpdfstring{\culine[\sectiontitleulinecolor]{\sectiontitle}}{\sectiontitle} }
\framesubtitle{Reaction cross section}

% 
\begin{itemize}
	\item Consider e.g. beam of protons incident on ${ }^{13} \mathrm{C}$ atoms:
	\item[] \begin{itemize}[itemsep=3mm]
	 \item \CDAlert[red]{Elastic scattering} is written as: \appear{$\hspace{3.5cm} { }^{13} \mathrm{C}(\mathrm{p}, \mathrm{p}){ }^{13} \mathrm{C}$} 

\item If the scattering is \CDAlert[blue]{inelastic}, we may write:
$
\hspace{1cm} \appear{{ }^{13} \mathrm{C} ( \mathrm{p}, \mathrm{p}^{\prime} ) {}^{13} \mathrm{C}^{*}}
$\\
\appear{since now the outgoing proton is $p^{\prime}$}\appear{, and the \\
target atom is excited}
\end{itemize}



\item Other possible reactions are: 
$$
 \begin{array}{ccc}
 \appear{(p, n) &: \qquad \qquad& { }^{13} \mathrm{C}(p, n) {}^{13} \mathrm{N}} \\ 
 \appear{\text {capture} &  : \qquad\qquad & { }^{13} \mathrm{C}(p, \gamma) {}^{14} \mathrm{N}} \\ 
 \appear{(p, \alpha) & : \qquad \qquad& { }^{13} \mathrm{C}(p, \alpha) { }^{10} \mathrm{B}}\end{array}
$$
\appear{Each reaction has its own partial cross section.} \appear{Then the total sum of the partial cross section will be}
\[
\sigma  \equal \appear{\sigma_{p, p}} ~\plus~ \appear{\sigma_{p, p^{\prime}}} ~\plus~ \appear{\sigma_{p, n}}  ~\plus~ \appear{\sigma_{p, \gamma} ~+~ \sigma_{p, \alpha} ~+~ \cdots}
\]
\end{itemize}

%\xdef\loc{south east}
%\begin{tikzpicture}[remember picture,overlay] % anchor=south
%    \node (A) [yshift=-0.25*\figYshift,xshift=\figXshift, anchor=\loc] at (current page.\loc) {\includegraphics[width=7cm]{Figures_booklet/SOM_12_Nuclear_physics_figures/SOM_Nuclear_physics_Figure_1.png}}; % 8cm max height
%\end{tikzpicture}

\endofslide 
%\endofsection
\end{frame}
%%%%%%%%%%%%%%%


%%%%%%%%%%%%%%%
\begin{frame}
\frametitle{\texorpdfstring{\culine[\sectiontitleulinecolor]{\sectiontitle}}{\sectiontitle} }
\framesubtitle{Intermediate Nucleus ('Compound Nucleus')}

\semismall

% 6.5 cm = midway, 10 cm = 3/4 
%\vspace{2.5mm}
\begin{minipage}{10.5cm}
\begin{itemize}[itemsep=1.7mm]
	\item Many of the \CDAlert[fuchsia]{low-energy nuclear reactions} proceed as two-stage processes 
	\item The projectile and target nuclide first form a temporary \CDAlert[red]{'compound nucleus'} \appear{which will be \Alert[red]{de-excited} via an emission of one or more particles}
	\item The de-excitation of the compound nucleus is a statistical process\appear{, which does not depend on how it was originally formed}

\item For instance, an incident $1 \mathrm{\;MeV}$ proton has a speed about $10^{7} \mathrm{\;m} / \mathrm{s}$\appear{, so for the proton to cross the nucleus it takes a time of:}
\[
\Delta t  \equal \fraca{\Delta x}{V} \equal \fraca{R}{V}  \equal \fraca{10^{-14} \mathrm{\;m}}{10^{7} \mathrm{\;m} / \mathrm{s}} \equal \appear{10^{-21} \mathrm{\;s}}
\]
\vspace{3mm}

\item A \CDAlert[blue]{lifetime} of a compound nucleus can be in the order of $10^{-16} \mathrm{~s}$\appear{, too short to be directly measured}\appear{, but long enough ($ 1000\;000 \times 10^{-21} \mathrm{\;s} $)}\appear{, to approximate that the decay is independent of how the compound nucleus was originally formed}
\end{itemize}
\end{minipage}
%\vspace{2.5mm}

\xdef\loc{south east}%
\begin{tikzpicture}[remember picture,overlay] % anchor=south
    \node (A) [yshift=-0.15*\figYshift,xshift=\figXshift, anchor=\loc] at (current page.\loc) {\includegraphics[height=8.5cm]{Figures_booklet/SOM_Tipler_Nuclear_physics_figures/SOM_Tipler_Nuclear_physics_figure_52}};% 8cm max height
\end{tikzpicture}


\endofslide 
%\endofsection
\end{frame}
%%%%%%%%%%%%%%%

%%%%%%%%%%%%%%%%
%\begin{frame}
%\frametitle{\texorpdfstring{\culine[\sectiontitleulinecolor]{\sectiontitle}}{\sectiontitle} }
%\framesubtitle{Reaction pathways: Entrance and exit channels}
%
%\seminormal 
%% 
%\begin{itemize}
%	\item Six nuclear reactions whose product is the compound nucleus ${ }_{7}{ }^{14} N^{*}$ and four ways in which ${ }_{7}^{14} N^{*}$ can decay if its excitation energy is $12 \mathrm{\;MeV}$. Other decay modes are possible if the excitation energy is greater, fewer are possible if this energy is less.
%\end{itemize}
%
%%\xdef\loc{south east}
%%\begin{tikzpicture}[remember picture,overlay] % anchor=south
%%    \node (A) [yshift=-0.25*\figYshift,xshift=\figXshift, anchor=\loc] at (current page.\loc) {\includegraphics[width=7cm]{Figures_booklet/SOM_12_Nuclear_physics_figures/SOM_Nuclear_physics_Figure_1.png}}; % 8cm max height
%%\end{tikzpicture}
%
%\endofslide 
%%\endofsection
%\end{frame}
%%%%%%%%%%%%%%%%

%%%%%%%%%%%%%%%
\begin{frame}
\frametitle{\texorpdfstring{\culine[\sectiontitleulinecolor]{\sectiontitle}}{\sectiontitle} }
\framesubtitle{Example on nuclear reactions}

\seminormal
% 
\begin{itemize}
%	\item Consider six nuclear reactions whose product is the compound nucleus ${ }^{14}_{\;7} N^{*}$ \appear{and four ways in which ${ }_{\;7}^{14} N^{*}$ can decay if its excitation energy is $12 \mathrm{\;MeV}$.} \appear{Other decay modes are possible if the excitation energy is greater}\appear{, fewer are possible if this energy is less}
	
	\item Calculate the $Q$-values for two example reactions\appear{, and determine whether they are exothermic or endothermic.}
\item[]
	
\item[] \Alert[red]{Reaction 1}: \vspace{-7mm}
$$
\begin{aligned}
 \appear{{ }_{\;5}^{10} \mathrm{B}}  ~\plus~ \appear{{ }_{0}^{1} \mathrm{n}} \quad &\rightarrow \quad \appear{{ }_{3}^{7} \mathrm{Li}}  ~\plus~ \appear{{ }_{2}^{4} \mathrm{He}} \\ 
 \\
\appear{Q}  \equal \appear{(10.012937 \;u + 1.008665 \;u -7.016003& \;u - 4.0026033 \;u)} \appear{\times 931.5 \mathrm{MeV} / \mathrm{u}} \appear{= 2.79 \mathrm{\;MeV} }
\end{aligned}
$$
\appear{Thus, the reaction is \CDAlert[red]{exothermic}}\appear{, energy is produced in the reaction}

\item[]
\item[] \Alert[blue]{Reaction 2}:\vspace{-7mm} 
$$
\begin{aligned}
 \appear{{ }_{3}^{7} \mathrm{Li}}  ~\plus~ \appear{{ }_{1}^{1} \mathrm{H}} \quad &\rightarrow \quad \appear{{ }_{4}^{7} \mathrm{Be}}  ~\plus~ \appear{{ }_{0}^{1} \mathrm{n}} \\
 \\
  \appear{Q}  \equal \appear{(7.016003 \;u+1.007825 \;u  - 7.0166928& \;u -1.008665 \;u)} \appear{\times 931.5 \mathrm{\;MeV} / \mathrm{u}}  \equal \appear{-1.64 \mathrm{\;MeV} }
 \end{aligned}
$$
\appear{Thus, the reaction is \CDAlert[blue]{endothermic}}\appear{, energy is consumed in the reaction}
\end{itemize}

%\xdef\loc{south east}
%\begin{tikzpicture}[remember picture,overlay] % anchor=south
%    \node (A) [yshift=-0.25*\figYshift,xshift=\figXshift, anchor=\loc] at (current page.\loc) {\includegraphics[width=7cm]{Figures_booklet/SOM_12_Nuclear_physics_figures/SOM_Nuclear_physics_Figure_1.png}}; % 8cm max height
%\end{tikzpicture}

\endofslide 
%\endofsection
\end{frame}
%%%%%%%%%%%%%%%

%%%%%%%%%%%%%%%
\begin{frame}
\frametitle{\texorpdfstring{\culine[\sectiontitleulinecolor]{\sectiontitle}}{\sectiontitle} }
\framesubtitle{Fission}

% 6.5 cm = midway, 10 cm = 3/4 
%\vspace{2.5mm}
\begin{minipage}{8.5cm}
	\begin{itemize}
	\item \Alert[red]{The heavy nuclei have less binding energy per nucleon} than the ones with intermediate mass number 
	
%	\item Lower internal energy means smaller mass and release of kinetic energy
\item Being deeper in the potential well\appear{, means smaller mass} \appear{and release of kinetic energy}
	

%	\item Quite similarly to spontaneous fission described earlier\appear{, also here the driving force is the reduction of \CDAlert[blue]{Coulombic repulsion}} \appear{(the protons are, on the average, further away from each other)}
\item Quite similarly to spontaneous fission described earlier\appear{, also here the driving force is the reduction of \CDAlert[blue]{Coulombic repulsion}} \appear{(the protons are, on the average, further away from each other)}

\item \CDAlert[fuchsia]{Particle-induced fission} is of practical use.%
\appear{\xdef\loc{south east}%
\begin{tikzpicture}[remember picture,overlay] % anchor=south
    \node (A) [yshift=-0.15*\figYshift,xshift=\figXshift, anchor=\loc] at (current page.\loc) {\includegraphics[width=3.7cm]{Figures_booklet/SOM_12_Nuclear_physics_figures/SOM_Nuclear_physics_Figure_28.png}}; % 8cm max height
\end{tikzpicture}}%
 \appear{ During our course, especially the fission reaction launched by colliding neutrons is discussed}
\end{itemize}
\end{minipage}
%\vspace{2.5mm}


\xdef\loc{north east}
\begin{tikzpicture}[remember picture,overlay] % anchor=south
    \node (A) [yshift=0.5*\figYshift,xshift=\figXshift, anchor=\loc] at (current page.\loc) {\includegraphics[width=4.9cm]{Figures_booklet/SOM_12_Nuclear_physics_figures/SOM_Nuclear_physics_Figure_27.png}}; % 8cm max height
\end{tikzpicture}

\endofslide 
%\endofsection
\end{frame}
%%%%%%%%%%%%%%%

%%%%%%%%%%%%%%%
\begin{frame}
\frametitle{\texorpdfstring{\culine[\sectiontitleulinecolor]{\sectiontitle}}{\sectiontitle} }
\framesubtitle{Fission}

\begin{itemize}
	\item As an example, when ${ }^{235} \mathrm{U}$ nucleus is excited by the capture of a neutron\appear{, it splits into two nuclei}\appear{, each having roughly one half of the initial mass}
	
\item \CDAlert[fuchsia]{Fission of ${ }^{235} \mathrm{U}$}: $\quad \appear{{ }^{235} \mathrm{U}}  ~\plus~  \appear{\mathrm{n}} \quad \rightarrow \quad { }\appear{^{92} \mathrm{Kr}}  ~\plus~ \appear{{ }^{142} \mathrm{Ba}}  ~\plus~ \appear{2 \mathrm{n}}  ~\plus~ \appear{179.4 \mathrm{\;MeV}$}

\item Coulombic repulsive force drives the \CDAlert[red]{fragments} apart, \appear{they will \CDAlert[red]{gain kinetic energy}.} \appear{Due to collision, this energy is \Alert[red]{converted into thermal energy}}

\item From the plot showing the binding energy \appear{(per nucleon)}\appear{, it is clear, that \CDAlert[blue]{fusion works best for $A \leq 20$}} \appear{and \CDAlert[fuchsia]{fission works best for $A \geq 200$}}

\item  If in fission, e.g. $A=240$ splits into two $A=120$ nuclei, \appear{then the released energy will be ${\Approx}1 \mathrm{\;MeV} / \mathrm{nucleon}$}\appear{, i.e. around $200 \mathrm{\;MeV} /\text{atom}$}
\end{itemize}

%\xdef\loc{south east}
%\begin{tikzpicture}[remember picture,overlay] % anchor=south
%    \node (A) [yshift=-0.25*\figYshift,xshift=\figXshift, anchor=\loc] at (current page.\loc) {\includegraphics[width=7cm]{Figures_booklet/SOM_12_Nuclear_physics_figures/SOM_Nuclear_physics_Figure_1.png}}; % 8cm max height
%\end{tikzpicture}

\endofslide 
%\endofsection
\end{frame}
%%%%%%%%%%%%%%%

%%%%%%%%%%%%%%%
\begin{frame}
\frametitle{\texorpdfstring{\culine[\sectiontitleulinecolor]{\sectiontitle}}{\sectiontitle} }
\framesubtitle{Fission}

\seminormal
\begin{itemize}[itemsep=1.5mm]
	\item In practise, again a compound nucleus will be first formed\appear{, which will eventually decay into two fragments}

\item An important example is \CDAlert[fuchsia]{neutron bombardment of uranium-235}\appear{, giving rise to a nucleus of an \Alert[red]{(excited) uranium-236}}

\item  uranium-236 will then decay into two fragments \appear{for which we have several different 'reaction channels' }

\item As the \CDAlert[blue]{end products} we may have e.g.: \appear{barium and krypton}\appear{, xenon and strontium}\appear{, tin and molybdenum}

\item And, as shown by the reaction schemes\appear{, along with each reaction there are some \CDAlert[blue]{extra neutrons} emerging}\appear{, plus $\gamma$ quanta (not written in equations)}
$$
\begin{aligned}
\appear{{ }_{\;92}^{235} \mathrm{U} ~ + ~ { }_{0}^{1} \mathrm{n} } \quad &\rightarrow \quad \appear{{ }_{56}^{141} \mathrm{Ba} ~ + ~ { }_{36}^{92} \mathrm{Kr} ~ + ~ 3\cdot{ }_{0}^{1} \mathrm{n} }\\
\appear{{ }_{\;92}^{235} \mathrm{U} ~ + ~ { }_{0}^{1} \mathrm{n}} \quad &\rightarrow \quad \appear{{ }_{\,54}^{140} \mathrm{Xe} ~ + ~ { }_{38}^{94} \mathrm{Sr} ~ + ~ 2\cdot{ }_{0}^{1} \mathrm{n}} \\
\appear{{ }^{235}_{\;92} \mathrm{U} ~ + ~ { }_{0}^{1} \mathrm{n}} \quad &\rightarrow \quad \appear{{ }_{\;50}^{132} \mathrm{Sn} ~ + ~ { }_{42}^{101} \mathrm{Mo} ~ + ~ 3\cdot{ }_{0}^{1} \mathrm{n}}
\end{aligned}
$$

\end{itemize}

%\xdef\loc{south east}
%\begin{tikzpicture}[remember picture,overlay] % anchor=south
%    \node (A) [yshift=-0.25*\figYshift,xshift=\figXshift, anchor=\loc] at (current page.\loc) {\includegraphics[width=7cm]{Figures_booklet/SOM_12_Nuclear_physics_figures/SOM_Nuclear_physics_Figure_1.png}}; % 8cm max height
%\end{tikzpicture}

\endofslide 
%\endofsection
\end{frame}
%%%%%%%%%%%%%%%

%%%%%%%%%%%%%%%
\begin{frame}
\frametitle{\texorpdfstring{\culine[\sectiontitleulinecolor]{\sectiontitle}}{\sectiontitle} }
\framesubtitle{Fission}

\begin{itemize}
\item[] \[
{ }_{\;92}^{235} U ~ + ~ { }_{0}^{1} n \quad \oldrightarrow \quad{ }_{\;92}^{236} U^{*} \quad \oldrightarrow \quad \ldots
\]
	\item Three example outcomes of the dissociation of uranium-236, given above\appear{, are only \Alert[fuchsia]{some examples} of the dominating channels} 
	
	\item %As shown by the Figure 8-19 (a), 
	In practise, there are \CDAlert[fuchsia]{many alternative outcomes!} %and the yield distribution of different product nuclei shows a bimodal shape 
	
	\item In general, uranium-236 fragments into two nuclei\appear{, plus some neutrons}\appear{, plus some energy}
	
	\item  The details of the process depend on neutron energy\appear{, which is \Alert[blue]{different for \CDAlert[blue]{thermal (slow)} neutrons}} \appear{(}$\appear{E_{kin}} \approx \appear{1.5 k_B T} \appear{\ll 0.1 \text{\,MeV}}$\appear{)} \appear{and \Alert[red]{faster neutrons}} \appear{($E_{kin} \Approx 1 \mathrm{\;MeV}$)}
	
	\item Note also, that only rarely a nucleus fissions into two exactly equally sized nuclei

%\item  the distribution will depend on neutron energy, being different for thermal (slow) neutrons and faster, $14 \mathrm{\;MeV}$, neutrons
\end{itemize}

%\xdef\loc{south east}
%\begin{tikzpicture}[remember picture,overlay] % anchor=south
%    \node (A) [yshift=-0.25*\figYshift,xshift=\figXshift, anchor=\loc] at (current page.\loc) {\includegraphics[width=7cm]{Figures_booklet/SOM_12_Nuclear_physics_figures/SOM_Nuclear_physics_Figure_1.png}}; % 8cm max height
%\end{tikzpicture}

\endofslide 
%\endofsection
\end{frame}
%%%%%%%%%%%%%%%

%%%%%%%%%%%%%%%
\begin{frame}
\frametitle{\texorpdfstring{\culine[\sectiontitleulinecolor]{\sectiontitle}}{\sectiontitle} }
\framesubtitle{Chain reaction}

\seminormal
\begin{itemize}[itemsep=1.9mm]
	\item Assume that each reaction \CDAlert[red]{requires a neutron to initiate} the reaction\appear{, and that \Alert[blue]{each reaction liberates $n$ product neutrons} } 
	\implies{ Each reaction will/can initiate $n$ other reactions } 
	
	\item If one reaction produces an amount of energy $E_{0}$\appear{, then the $j$:th generation of fission could produce an energy of }
$ \qquad
\appear{E_{j}}  \equal \appear{E_{0}} \appear{n^{j}}
$

\item In 1942, a group led by \Alert[red]{Fermi} created the first chain reaction 

\item It is possible to show that this kind of chain reaction requires the fissioning material to have a mass larger than so-called  \CDAlert[fuchsia]{critical mass}

\item Many neutrons can emerge from each reaction \appear{$(n>1)$}

\item But not all neutrons initiate further reactions \appear{(losses etc.)}\appear{, in practice it is better to replace the equation above by}
\[
E_{j}  \equal \appear{E_{0} k^{j} \,,} 
\]
\appear{\CDAlert[fuchsia]{if $k>1$ we will have a chain reaction}}\appear{, if $k<1$, the reaction will die out}
\end{itemize}

%\xdef\loc{south east}
%\begin{tikzpicture}[remember picture,overlay] % anchor=south
%    \node (A) [yshift=-0.25*\figYshift,xshift=\figXshift, anchor=\loc] at (current page.\loc) {\includegraphics[width=7cm]{Figures_booklet/SOM_12_Nuclear_physics_figures/SOM_Nuclear_physics_Figure_1.png}}; % 8cm max height
%\end{tikzpicture}

\endofslide 
%\endofsection
\end{frame}
%%%%%%%%%%%%%%%

%%%%%%%%%%%%%%%
\begin{frame}
\frametitle{\texorpdfstring{\culine[\sectiontitleulinecolor]{\sectiontitle}}{\sectiontitle} }
\framesubtitle{Nuclear reactions}

% on nuclear reactions
\begin{itemize}
	\item In a nuclear reactor\appear{, based on fission reaction}\appear{, a fuel assembly to provide $k>1$ is used }
	
	\item \CDAlert[fuchsia]{Control rods }are used to keep $k$ at $\Approx 1$\appear{, by slowing down} \appear{or capturing the neutrons and thereby decreasing $k$}
	
	\item boron and cadmium are commonly used materials (recall earlier lecture)

	\item A fraction of neutrons $(\sim 1 \%)$ from the decay have a significant time% lag (${\approx}1\,s$ delay),
\end{itemize}

% 6.5 cm = midway, 10 cm = 3/4 
%\vspace{2.5mm}
\begin{minipage}{5.15cm}
\begin{itemize}
	\item[] lag \appear{(${\Approx}1\,s$ delay)}\appear{, which slows down the overall chain reaction }
\implies{ Enables to control of the \\
	      process and avoiding $k$\\
	      to increase above unity }
\end{itemize}
\end{minipage}
%\vspace{2.5mm}


\xdef\loc{south east}
\begin{tikzpicture}[remember picture,overlay] % anchor=south
    \node (A) [yshift=-0.25*\figYshift,xshift=\figXshift, anchor=\loc] at (current page.\loc) {\includegraphics[width=8.5cm]{Figures_booklet/SOM_12_Nuclear_physics_figures/SOM_Nuclear_physics_Figure_29.png}}; % 8cm max height
\end{tikzpicture}

%\endofslide 
\endofsection
\end{frame}
%%%%%%%%%%%%%%%


% fusion
%%%%%%%%%%%%%%%
\begin{frame}
\frametitle{\texorpdfstring{\culine[\sectiontitleulinecolor]{\sectiontitle}}{\sectiontitle} }
\framesubtitle{Fusion}

% 
\begin{itemize}
	\item Light nuclei are less tightly bound than the intermediate ones. \appear{When light nuclei are merged into heavier ones}\appear{, total mass will decrease and kinetic energy will increase }
	
	\item This process is called  \CDAlert[fuchsia]{fusion}\appear{, and a typical reaction is of form}
\[
\appear{{ }_{1}^{2} \mathrm{H}} \appear{ ~+~ { }_{1}^{1} \mathrm{H}} \quad \oldrightarrow \quad { }_{2}^{3} \mathrm{He} ~+~ \gamma \appear{\,, \qquad Q=5.48 \mathrm{\;MeV}}
\]
% more stuff into equation
\end{itemize}

% 6.5 cm = midway, 10 cm = 3/4 
%\vspace{2.5mm}
\begin{minipage}{7.5cm}
\begin{itemize}
	\item Protons, having \CDAlert[blue]{Coulombic repulsion}\appear{, are packed into same nucleus} \appear{(requires an energy in the order of $1 \mathrm{\;MeV}$)} 
	
	\item But in the same time more nucleons get close to each other within the range of strong force\appear{, ${\sim}2 \times 10^{-15} \mathrm{\;m}$} 
	\implies{ \Alert[red]{Number of bonds between nucleons \\
			increases}}
\end{itemize}
\end{minipage}
%\vspace{2.5mm}


\xdef\loc{south east}
\begin{tikzpicture}[remember picture,overlay] % anchor=south
    \node (A) [yshift=-0.25*\figYshift,xshift=\figXshift, anchor=\loc] at (current page.\loc) {\includegraphics[width=5.5cm]{Figures_booklet/SOM_12_Nuclear_physics_figures/SOM_Nuclear_physics_Figure_30.png}}; % 8cm max height
\end{tikzpicture}

\endofslide 
%\endofsection
\end{frame}
%%%%%%%%%%%%%%%

%%%%%%%%%%%%%%%
\begin{frame}
\frametitle{\texorpdfstring{\culine[\sectiontitleulinecolor]{\sectiontitle}}{\sectiontitle} }
\framesubtitle{Fusion}

\seminormal
\begin{itemize}
	\item Fusion reactions can take place via thermal activation at \Alert[red]{very high temperatures}

\item A sufficient fraction of nuclei have this energy \appear{when $k_{B} T \Approx 10 \mathrm{\;keV}$}\appear{, thus $T \Approx 10^{8} \mathrm{\;K}$}

\item These high temperatures can be found only \CDAlert[red]{within the stars}\appear{, where matter is in ionized plasma phase} \appear{and gravitational forces hold the material together}

\item For the fusion to happen\appear{, the \CDAlert[red]{hydrogen-plasma} will need to be held together in a particle number density of matter $n$ for a residence time of $\tau$}\appear{, and together they need to obey the condition:}
\[
n \appear{\tau} \appear{>} \appear{10^{20}} \; \frac{\text {particles} \cdot s}{m^{3}} \qquad \appear{(\CDAlert[fuchsia]{Lawson~criterion})}
\]

\item Here a certain residence time $\tau$ \appear{is needed because reaction takes some time}

\end{itemize}

%\xdef\loc{south east}
%\begin{tikzpicture}[remember picture,overlay] % anchor=south
%    \node (A) [yshift=-0.25*\figYshift,xshift=\figXshift, anchor=\loc] at (current page.\loc) {\includegraphics[width=7cm]{Figures_booklet/SOM_12_Nuclear_physics_figures/SOM_Nuclear_physics_Figure_1.png}}; % 8cm max height
%\end{tikzpicture}

\endofslide 
%\endofsection
\end{frame}
%%%%%%%%%%%%%%%

%%%%%%%%%%%%%%%
\begin{frame}
\frametitle{\texorpdfstring{\culine[\sectiontitleulinecolor]{\sectiontitle}}{\sectiontitle} }
\framesubtitle{Proton–proton cycle}

\seminormal
\begin{itemize}
	\item Hydrogen is by far the most abundant element in the stars\appear{, and a series of fusion reactions known as the \CDAlert[fuchsia]{proton–proton cycle}} \appear{"burns" hydrogen towards helium-4}
$$
\begin{aligned}
& \appear{\Alert[red]{Step 1}:}\quad & \appear{{ }_{1}^{1} H}  ~\plus~ \appear{{ }_{1}^{1} H} \quad &\rightarrow \quad \appear{{ }_{1}^{2} H}  ~\plus~ \appear{{ }_{+1}^{\,0} \beta ~+~ \nu} \appear{&Q&=0.42 \mathrm{\;MeV}} \\
& \appear{\Alert[blue]{Step 2}:}\quad & \appear{{ }_{1}^{2} \mathrm{H} ~+~ { }_{1}^{1} \mathrm{H}} \quad &\rightarrow \appear{\quad { }_{2}^{3} \mathrm{He}}  \appear{&Q&=5.48 \mathrm{\;MeV}} \\
& \appear{\Alert[green]{Step 3}:} \qquad & \appear{{ }_{2}^{3} \mathrm{He}}  ~\plus~ \appear{{ }_{2}^{3} \mathrm{He}} \quad &\rightarrow \appear{\quad { }_{2}^{4} \mathrm{He}}  ~\plus~ \appear{{ }_{1}^{1} \mathrm{H} ~+~ { }_{1}^{1} \mathrm{H}} &\quad \appear{Q&=12.9 \mathrm{\;MeV}}
\end{aligned}
$$

\item The cycle begins with two protons forming a \CDAlert[red]{deuteron} (\Alert[red]{Step 1}). \appear{For this fusion, a proton needs to become a neutron}\appear{, therefore the positron and the neutrino will be formed} 

\item Because this reaction is mediated by the \Alert[red]{weak force}\appear{, it takes some time }

\item The formed deuteron combines with another proton (\Alert[blue]{Step 2})\appear{, and forms \CDAlert[blue]{helium-3}}\appear{, which, with another helium-3 reacts to form \CDAlert[green]{helium-4} (\Alert[green]{Step 3})}

\item Note, two of the steps $1 \;\&\; 2$ are needed \appear{for one step $3$ to take place}
\end{itemize}

%\xdef\loc{south east}
%\begin{tikzpicture}[remember picture,overlay] % anchor=south
%    \node (A) [yshift=-0.25*\figYshift,xshift=\figXshift, anchor=\loc] at (current page.\loc) {\includegraphics[width=7cm]{Figures_booklet/SOM_12_Nuclear_physics_figures/SOM_Nuclear_physics_Figure_1.png}}; % 8cm max height
%\end{tikzpicture}

\endofslide 
%\endofsection
\end{frame}
%%%%%%%%%%%%%%%

%%%%%%%%%%%%%%%
\begin{frame}
\frametitle{\texorpdfstring{\culine[\sectiontitleulinecolor]{\sectiontitle}}{\sectiontitle} }
\framesubtitle{Proton–proton cycle}

\seminormal
% 6.5 cm = midway, 10 cm = 3/4 
%\vspace{2.5mm}
\begin{minipage}{7.5cm}
	\begin{itemize}
	\item In total \CDAlert[fuchsia]{four protons turn into one He-4} (plus some extra positrons and neutrinos)

\item Also $24.7 \mathrm{\;MeV}$ of energy is released:
$$\def\arraystretch{1.15}
\begin{array}{crr} 
2 \times \text { step } 1 & \qquad & 2 \times 0.42 \mathrm{\;MeV} \\ 
+2 \times \text { step } 2 & \qquad & +2 \times 5.48 \mathrm{\;MeV} \\ 
+1 \times \text { step } 3  & \qquad & +1 \times 12.9 \mathrm{\;MeV} \\ 
\hline
&  & =\phantom{+2\;} 24.7 \mathrm{\;MeV}
\end{array}
$$

\item In total:
\[
4 \cdot { }_{1}^{1} H \quad \rightarrow \quad \appear{{ }_{2}^{4} \mathrm{He}} \appear{ ~+~ 2 \beta^{+} ~+~ 2 \mathrm{\nu} ~+~ 24.7 \mathrm{\;MeV}}
\]

\item This \CDAlert[fuchsia]{proton–proton cycle} is the primary energy source in our Sun. It is estimated that the reaction is occurring at the rate of $5.64 \times 10^{11} \mathrm{\;kg}$ of hydrogen per second fusing into helium producing $3.7 \times 10^{25} \;W$
 \end{itemize}
\end{minipage}
%\vspace{2.5mm}


\xdef\loc{east}
\begin{tikzpicture}[remember picture,overlay] % anchor=south
    \node (A) [yshift=-0.0*\figYshift,xshift=\figXshift, anchor=\loc] at (current page.\loc) {\includegraphics[height=8cm]{Figures_booklet/SOM_12_Nuclear_physics_figures/SOM_Nuclear_physics_Figure_31.png}}; % 8cm max height
\end{tikzpicture}

\endofslide 
%\endofsection
\end{frame}
%%%%%%%%%%%%%%%

%%%%%%%%%%%%%%%
\begin{frame}
\frametitle{\texorpdfstring{\culine[\sectiontitleulinecolor]{\sectiontitle}}{\sectiontitle} }
\framesubtitle{Carbon cycle}

\semismall
% 
\begin{itemize}
	\item Proton–proton cycle is assumed to be the dominating reaction occurring in our Sun\appear{, but in other \CDAlert[blue]{stars of higher temperature and higher helium} \CDAlert[blue]{concentration}}\appear{, two helium-4 nuclei can form beryllium-8} 
	\[
	 { }_{2}^{4} \mathrm{He} ~+~ { }_{2}^{4} \mathrm{He} \quad \rightarrow \quad \appear{{ }_{4}^{8} \mathrm{Be} }
	 \]

\item ${ }_{4}^{8} Be$ is unstable and decays back to two heliums (in ${\sim} 10^{-16} \mathrm{\;s} $)\appear{, but if another ${ }_{2}^{4} \mathrm{He}$ fusions with Be before its decay}\appear{, a tightly bound ${ }_{\;6}^{12} C$ can be formed:}
\end{itemize}

% 6.5 cm = midway, 10 cm = 3/4 
\vspace{2.5mm}
\begin{minipage}{7.7cm}
\begin{itemize}
	\item[] \[
{ }_{4}^{8} \mathrm{Be}  ~\plus~ \appear{{ }_{2}^{4}\mathrm{He}} \quad \rightarrow \quad \appear{{ }_{\;6}^{12} C}
\]

\item This C nucleus can begin the \CDAlert[fuchsia]{carbon cycle}\appear{, which is a rapid way to "burn" hydrogen}
\appear{\xdef\loc{south east}
\begin{tikzpicture}[remember picture,overlay] % anchor=south
    \node (A) [yshift=-0.25*\figYshift,xshift=\figXshift, anchor=\loc] at (current page.\loc) {\semismall $\begin{aligned}
&1:& \quad  {\color{blue}{ }_{\;6}^{12} \mathrm{C}} + { }_{1}^{1} \mathrm{H} \quad &\oldrightarrow \quad { }_{\;7}^{13} \mathrm{N}  \\
&2:& \quad \quad { }_{\;7}^{13} \mathrm{N} \quad \quad &\oldrightarrow \quad { }_{\;6}^{13} \mathrm{C}+{ }_{+1}^{0} \beta+\nu \\
&3:& \quad { }_{\;6}^{13} \mathrm{C}+{ }_{1}^{1} \mathrm{H} \quad &\oldrightarrow \quad { }_{\;7}^{14} \mathrm{N} \\
&4:& \quad { }_{\;7}^{14} \mathrm{N}+{ }_{1}^{1} \mathrm{H} \quad &\oldrightarrow \quad { }_{\;8}^{15} \mathrm{O} \\
&5:& \quad  { }_{\;8}^{15} \mathrm{O} \quad \quad &\oldrightarrow  \quad { }_{\;7}^{15}{N} + { }_{+1}^{\;0}\beta+\nu \\
&6:& \quad { }_{\;7}^{15} \mathrm{N}+{ }_{1}^{1} \mathrm{H} \quad &\oldrightarrow \quad {\color{blue}{ }_{\;6}^{12} \mathrm{C}} +{ }_{2}^{4} \mathrm{He}
\end{aligned}$}; % 8cm max height
\end{tikzpicture}}

\item Note, that the \CDAlert[blue]{carbon nucleus} merely acts as a \CDAlert[blue]{catalyst} in this reaction scheme. \appear{Note also that there are heavier nuclei} \appear{such as ${ }_{\;8}^{15} \mathrm{O}$} \appear{and ${ }^{15} N$ in the cycle}\appear{, but they decay normally}
\end{itemize}
\end{minipage}
%\vspace{2.5mm}



%$$
%\begin{aligned}
%&1:{ }_{6}^{12} \mathrm{C}+{ }_{1}^{1} \mathrm{H} \rightarrow { }_{7}^{13} \mathrm{~N}  
%&2&:{ }_{7}^{13} \mathrm{~N} \rightarrow{ }_{6}^{13} \mathrm{C}+{ }_{+1}^{0} \beta+\nu \\
%&3: { }_{6}^{13} \mathrm{C}+{ }_{1}^{1} \mathrm{H} \rightarrow{ }_{7}^{14} \mathrm{~N} 
%&4&:{ }_{7}^{14} \mathrm{~N}+{ }_{1}^{1} \mathrm{H} \rightarrow{ }_{8}^{15} \mathrm{O} \\
%&5:{ }_{8}^{15} \mathrm{O} \rightarrow{ }_{7}^{15} \mathrm{~N}+{ }_{+1}^{0} \beta+\nu 
%&\qquad 6&:{ }_{7}^{15} \mathrm{~N}+{ }_{1}^{1} \mathrm{H} \rightarrow{ }_{6}^{12} \mathrm{C}+{ }_{2}^{4} \mathrm{He}
%\end{aligned}
%$$

%\xdef\loc{south east}
%\begin{tikzpicture}[remember picture,overlay] % anchor=south
%    \node (A) [yshift=-0.25*\figYshift,xshift=\figXshift, anchor=\loc] at (current page.\loc) {\includegraphics[width=7cm]{Figures_booklet/SOM_12_Nuclear_physics_figures/SOM_Nuclear_physics_Figure_1.png}}; % 8cm max height
%\end{tikzpicture}

\endofslide 
%\endofsection
\end{frame}
%%%%%%%%%%%%%%%

%%%%%%%%%%%%%%%
\begin{frame}
\frametitle{\texorpdfstring{\culine[\sectiontitleulinecolor]{\sectiontitle}}{\sectiontitle} }
\framesubtitle{Carbon cycle}

\begin{itemize}
	\item Note, that as the outcome of the cycle\appear{, four protons produce one helium-4 and two positrons and two neutrinos} 
	
	\item But this reaction is \CDAlert[fuchsia]{much faster than the proton–proton cycle}\appear{, because at no point does it require a proton to be converted into a neutron} \appear{before the fusion reaction can happen}
	
%$$
%\begin{aligned}
%&{ }_{1}^{1} \mathrm{H}+{ }_{6}^{12} \mathrm{C} \rightarrow{ }_{7}^{13} \mathrm{~N} \\
%&{ }^{13} \mathrm{~N} \rightarrow{ }_{6}^{13} \mathrm{C}+\mathrm{e}^{+}+\nu \\
%&{ }_{1}^{1} \mathrm{H}+{ }_{6}^{13} \mathrm{C} \rightarrow{ }_{7}^{14} \mathrm{~N} \\
%&{ }_{1}^{1} \mathrm{H}+{ }_{7}^{14} \mathrm{~N} \rightarrow{ }_{8}^{15} \mathrm{O} \\
%&{ }^{15}{ }_{8}^{15} \mathrm{O} \rightarrow{ }_{7}^{15} \mathrm{~N}+\mathrm{e}^{+}+\nu \\
%&{ }_{1}^{1} \mathrm{H}+{ }_{7}^{15} \mathrm{~N} \rightarrow{ }_{6}^{12} \mathrm{C}+{ }_{2}^{4} \mathrm{He} \\
%&4{ }_{1}^{1} \mathrm{H} \rightarrow{ }_{2}^{4} \mathrm{He}+2 \mathrm{e}^{+}+2 \nu+26.7 \mathrm{\;MeV}
%\end{aligned}
%$$
	
\item The energy produced by one cycle is $26.7 \mathrm{\;MeV}$ \appear{i.e. almost equal to the energy released in the proton–proton cycle:}
\[
4 \cdot { }_{1}^{1} \mathrm{H} \quad \rightarrow \quad \appear{{ }_{2}^{4} \mathrm{He}} \appear{ ~+~ 2 \mathrm{e}^{+} ~+~ 2 \nu}
\]

\item In our Sun, proton–proton cycle dominates 

\item \CDAlert[fuchsia]{In older and hotter stars}\appear{, carbon cycle as well as also other types of fusion reactions are known to occur}

\end{itemize}

%\xdef\loc{south east}
%\begin{tikzpicture}[remember picture,overlay] % anchor=south
%    \node (A) [yshift=-0.25*\figYshift,xshift=\figXshift, anchor=\loc] at (current page.\loc) {\includegraphics[width=7cm]{Figures_booklet/SOM_12_Nuclear_physics_figures/SOM_Nuclear_physics_Figure_1.png}}; % 8cm max height
%\end{tikzpicture}

\endofslide 
%\endofsection
\end{frame}
%%%%%%%%%%%%%%%

%%%%%%%%%%%%%%%
\begin{frame}
\frametitle{\texorpdfstring{\culine[\sectiontitleulinecolor]{\sectiontitle}}{\sectiontitle} }
\framesubtitle{Nucleosynthesis: Origin of elements}

%
\begin{itemize}
	\item In the universe\appear{, \CDAlert[blue]{hydrogen is most abundant element}}\appear{, \CDAlert[blue]{then helium}} 
	
	\item After the sudden drop\appear{, corresponding to lithium, beryllium and boron}\appear{, the \CDAlert[red]{abundance follows a regularly decreasing trend} }
	
	\item The same pattern covers not only samples taken from various depths from Earth's crust\appear{, but also from the meteorites}
	
\item The abundance more or less level off at around $Z=35$ \appear{(or $A=80)$}\appear{, but still it does have some pronounced maxima}

\item As for isotopic composition\appear{, lighter elements are richer with the stable isotopes with lower neutron content}\appear{, while the heavier nuclei are richer in isotopes with higher neutron content} 

\item Interestingly, (stable) isotopes with $A=5$ and $A=8$ \appear{are missing from nature}
\end{itemize}

%\xdef\loc{south east}
%\begin{tikzpicture}[remember picture,overlay] % anchor=south
%    \node (A) [yshift=-0.25*\figYshift,xshift=\figXshift, anchor=\loc] at (current page.\loc) {\includegraphics[width=7cm]{Figures_booklet/SOM_12_Nuclear_physics_figures/SOM_Nuclear_physics_Figure_1.png}}; % 8cm max height
%\end{tikzpicture}

\endofslide 
%\endofsection
\end{frame}
%%%%%%%%%%%%%%%

%%%%%%%%%%%%%%%
\begin{frame}
\frametitle{\texorpdfstring{\culine[\sectiontitleulinecolor]{\sectiontitle}}{\sectiontitle} }
\framesubtitle{Nucleosynthesis: Origin of elements}

\semismall
\begin{itemize}[itemsep=1.4mm]
	\item It seems that, at least in our galaxy\appear{, the elements were all formed by approximately the same process}
	
	\item At an early time in the universe\appear{, there was a very large mass of gaseous hydrogen.} \appear{Due to statistical fluctuations and gravitational interaction}\appear{, some \CDAlert[fuchsia]{hydrogen condensed into clusters}}\appear{, and stars with high density $10^{5} \mathrm{\;kg} / \mathrm{m}^{3}$}
	
	\item \CDAlert[fuchsia]{Temperatures in stars/clusters increased} \appear{and enabled proton–proton cycle}\appear{, After that the helium-4 lead the reactions to larger nuclei }
	
	\item Note, ${ }_{4}^{8} \mathrm{Be}$ is not stable\appear{, but fusion with ${ }_{1}^{2} \mathrm{H}$ or ${ }_{2}^{3} \mathrm{He}$ may lead to ${ }_{\;5}^{10} B$, ${ }_{\;6}^{11} C$ etc.}
	
	\item The schematic presentation of a reaction where the beryllium captures a helium-4 \appear{and leads to so called \CDAlert[blue]{'helium burning' process}:}
	$$
\begin{aligned}
&\appear{{ }_{2}^{4} \mathrm{He}}  ~\plus~ \appear{{ }_{2}^{4} \mathrm{He}} \quad &\rightarrow& \qquad \appear{{ }_{4}^{8}\mathrm{Be}} \qquad &\appear{(-0.0918 \mathrm{\;MeV})} \\
&\appear{{ }_{4}^{8} \mathrm{Be} ~+~ { }_{2}^{4} \mathrm{He}} \quad &\rightarrow& \quad \appear{{ }_{\;6}^{12} \mathrm{C} ~+~ 2 \gamma \qquad &(+7.367 \mathrm{\;MeV})}
\end{aligned}
$$

\item Ordinary fusion reactions reach to $A \Approx 60$. \appear{Beyond that, \CDAlert[red]{neutron capture} followed by \CDAlert[red]{$\beta^{-}$ decay} leads to higher $Z$} \appear{(\Alert[red]{supernovas}, merging neutron stars, etc.)}
\end{itemize}

%\xdef\loc{south east}
%\begin{tikzpicture}[remember picture,overlay] % anchor=south
%    \node (A) [yshift=-0.25*\figYshift,xshift=\figXshift, anchor=\loc] at (current page.\loc) {\includegraphics[width=7cm]{Figures_booklet/SOM_12_Nuclear_physics_figures/SOM_Nuclear_physics_Figure_1.png}}; % 8cm max height
%\end{tikzpicture}

\endofslide 
%\endofsection
\end{frame}
%%%%%%%%%%%%%%%

%%%%%%%%%%%%%%%
\begin{frame}
\frametitle{\texorpdfstring{\culine[\sectiontitleulinecolor]{\sectiontitle}}{\sectiontitle} }
\framesubtitle{Inducing the fusion reaction}

%
\begin{itemize}
	\item Fusion cannot occur spontaneously. \appear{In order to \Alert[red]{surpass the Coulombic repulsion} (nuclei are positively charged)}\appear{, nuclei need to have large kinetic energies (= \CDAlert[red]{high temperature})} 
\end{itemize}
\appear{\xdef\loc{south east}%
\begin{tikzpicture}[remember picture,overlay] % anchor=south
    \node (A) [yshift=-0.15*\figYshift,xshift=\figXshift, anchor=\loc] at (current page.\loc) {\includegraphics[height=6.9cm]{Figures_booklet/SOM_12_Nuclear_physics_figures/SOM_Nuclear_physics_Figure_32.png}}; % 8cm max height
\end{tikzpicture}}
	
% 6.5 cm = midway, 10 cm = 3/4 
\vspace{-0.3mm}
\begin{minipage}{7.5cm}
\begin{itemize}[itemsep=2.2mm]
	\item Once the nuclei are close to each other\appear{, within the range of the strong force}\appear{, the nuclei can then "fall into the potential pit" (see Fig.)}

\item It is obvious that fusion requires high density and high temperatures 

\item Once reactions start happening\appear{, large amount of energy will be released}\appear{, as can be deduced form the steep left edge of the binding energy curve}
\end{itemize}
\end{minipage}
%\vspace{2.5mm}

\endofslide 
%\endofsection
\end{frame}
%%%%%%%%%%%%%%%

%%%%%%%%%%%%%%%
\begin{frame}
\frametitle{\texorpdfstring{\culine[\sectiontitleulinecolor]{\sectiontitle}}{\sectiontitle} }
\framesubtitle{Fusion reactors}

% 
\begin{itemize}
	\item A fusion reaction, which has a large cross section \appear{and which liberates a great amount of energy}\appear{, is the reaction $D+T$}\appear{, i.e. between \CDAlert[red]{deuterium} $(D)$ and \CDAlert[blue]{tritium} $(T)$} 
	
	\item However, since tritium is not readily available\appear{, the \CDAlert[red]{fusion of two} \CDAlert[red]{deuterium nuclei $(D+D)$} is of more practical importance }
	
	\item Also, it has the advantage that only one class of nuclei ($D$) is used 
	
	\item Furthermore, $D+D$ fusion produces also tritium needed for the second reaction:
	$$
\appear{{ }_{\,1}^{2} \mathrm{D}}  ~\plus~ \appear{{ }_{\,1}^{2} \mathrm{D}} \quad \rightarrow \quad 
\LC[3] \begin{array}{ll} \appear{{ }^{3} \mathrm{T}} ~\plus~ \appear{{ }_{1}^{1} \mathrm{H}} & \appear{\qquad Q=4.0 \mathrm{\;MeV}} \\ \appear{{ }_{2}^{3} \mathrm{He}}  ~\plus~ \appear{{ }_{0}^{1} \mathrm{n}} & \qquad \appear{Q=3.3 \mathrm{\;MeV}} 
\end{array}
$$
and 
\[
{ }_{\,1}^{2} D  ~\plus~ \appear{{ }_{\,1}^{3} \mathrm{T}} \quad \rightarrow \quad \appear{{ }_{2}^{4} \mathrm{He}}  ~\plus~ \appear{{ }_{0}^{1} \mathrm{n}} \qquad \appear{Q=17.6 \mathrm{\;MeV}}
\]

\end{itemize}

%\xdef\loc{south east}
%\begin{tikzpicture}[remember picture,overlay] % anchor=south
%    \node (A) [yshift=-0.25*\figYshift,xshift=\figXshift, anchor=\loc] at (current page.\loc) {\includegraphics[width=7cm]{Figures_booklet/SOM_12_Nuclear_physics_figures/SOM_Nuclear_physics_Figure_1.png}}; % 8cm max height
%\end{tikzpicture}

\endofslide 
%\endofsection
\end{frame}
%%%%%%%%%%%%%%%

%%%%%%%%%%%%%%%
\begin{frame}
\frametitle{\texorpdfstring{\culine[\sectiontitleulinecolor]{\sectiontitle}}{\sectiontitle} }
\framesubtitle{Fusion reactors}

\semismall
% 6.5 cm = midway, 10 cm = 3/4 
%\vspace{2.5mm}
\begin{minipage}{8.5cm}
\begin{itemize}[itemsep=1.6mm]
	\item One needs high energy for fusion. %, to bring the nuclei close to each other, against the Coulombic potential barrier (shown in Figure 32). 
	%\item 
	\appear{The calculated temperature required to bring two deuterons within a distance of ${\sim} 2 \mathrm{\;fm}$ from each other is $T \sim 6 \times 10{ }^{9} \mathrm{\;K}$} 
	
	\item There is, naturally, a velocity distribution\appear{, so some nuclei have a higher velocity than the average.} \appear{Also tunneling helps.} \appear{All in all,  temperatures of $10^{7}–10^{8} \;K$} \appear{may be sufficient to have the reaction occurring (as is the case in the Sun)}

\item In the real \CDAlert[fuchsia]{fusion reactors}\appear{, the \Alert[blue]{challenge is to confine the plasma}} \appear{exhibiting a sufficient temperature for the reactions to happen}

\item Two promising techniques are available:\\
\appear{\CDAlert[red]{Magnetic confinement}} $\rightarrow$ \appear{toroidal magnetic field}\\
\appear{\CDAlert[blue]{Inertial confinement}} $\rightarrow$ \appear{synchronized blows of the sample pellet from all sides}\appear{, via laser beams illuminating the pellet}

%\item $$
%\begin{gathered}
%\frac{1}{4 \pi \varepsilon_{0}} \frac{q_{1} q_{2}}{r}=\frac{3}{2} k_{B} T \\
%\Rightarrow T \cong 6 \times 10^{9} K
%\end{gathered}
%$$
\end{itemize}

\end{minipage}
%\vspace{2.5mm}

\xdef\loc{south east}
\begin{tikzpicture}[remember picture,overlay] % anchor=south
    \node (A) [yshift=-0.25*\figYshift,xshift=\figXshift, anchor=\loc] at (current page.\loc) {\includegraphics[width=5cm]{Figures_booklet/SOM_12_Nuclear_physics_figures/SOM_Nuclear_physics_Figure_33.png}}; % 8cm max height
\end{tikzpicture}

\endofslide 
%\endofsection
\end{frame}
%%%%%%%%%%%%%%%

%%%%%%%%%%%%%%%
\begin{frame}
\frametitle{\texorpdfstring{\culine[\sectiontitleulinecolor]{\sectiontitle}}{\sectiontitle} }
\framesubtitle{Comparison between fission and fusion}

%
\begin{itemize}
	\item[] \begin{itemize}
	\item[\textbf{\color{red} +}] Both techniques, fission and fusion, are based on strong force
\item[\textbf{\color{red} +}] Only a small amount of fuel is needed \appear{(compared to fossil fuels)}
\item[\textbf{\color{red} +}] No waste $\mathrm{CO}_{2}$ emitted $\rightarrow$ \appear{atmospheric sustainability}
\item[\textbf{\color{blue} -}] Reaction products are radioactive $\rightarrow$ \appear{other environmental disadvantages}
\appear{($\oldrightarrow$ final disposal??)}
\end{itemize}

\item \CDAlert[blue]{Fission}:
\begin{itemize}
	\item Fuels: \appear{uranium and thorium, need for expensive mining processes}
\item Process products: \appear{radioactive waste with long half-life} \appear{(e.g. $T_{1/2}(\mathrm{Pu-239}) =24\;100$ years)}
\end{itemize}

\item \CDAlert[red]{Fusion}:
\begin{itemize}
	\item Fuel: \appear{deuterium can be distilled from water}\appear{, is abundant}
\item Process products: \appear{tritium and isotopes of helium}\appear{, tritium is radioactive}\appear{, but its half-life is short (12.3 years)}
\end{itemize}

\end{itemize}

%\xdef\loc{south east}
%\begin{tikzpicture}[remember picture,overlay] % anchor=south
%    \node (A) [yshift=-0.25*\figYshift,xshift=\figXshift, anchor=\loc] at (current page.\loc) {\includegraphics[width=7cm]{Figures_booklet/SOM_12_Nuclear_physics_figures/SOM_Nuclear_physics_Figure_1.png}}; % 8cm max height
%\end{tikzpicture}

\endofslide 
%\endofsection
\end{frame}
%%%%%%%%%%%%%%%

%%%%%%%%%%%%%%%
\begin{frame}
\frametitle{\texorpdfstring{\culine[\sectiontitleulinecolor]{\sectiontitle}}{\sectiontitle} }
\framesubtitle{On the energy yields}

\seminormal
\begin{itemize}
	\item Energy released in a single fusion reaction is much less than that released in fission. \appear{But normalized against the mass}\appear{, for the deuteron–deuteron fusion reaction the energy is about $2 \times 10^{14} \;J$ per kg of fuel}
	\implies{ Over double the value for uranium fission }

\item Because deuterium is abundant \appear{(about 1 D per 7000 hydrogen atoms)}\appear{, and the cost of extracting D from water is low (a fraction of €/gram)}\appear{, fusion is considered to be the \CDAlert[fuchsia]{'promising power source of the future'}}

\item However, as has been empirically seen \appear{(e.g. ITER-reactors)}\appear{, it is still \Alert[fuchsia]{technically very challenging}} 

\item Note also the timeline and what has been stated about fusion power:
\begin{itemize}
	\item "A promising power source to be available within 20 years..."\\
	 \appear{– statement from year $1985 !$}
\item "Still a promising power source to be available within 20 years..."\\
	  \appear{– statement in year $2020 !$}
\end{itemize}

\end{itemize}

%\xdef\loc{south east}
%\begin{tikzpicture}[remember picture,overlay] % anchor=south
%    \node (A) [yshift=-0.25*\figYshift,xshift=\figXshift, anchor=\loc] at (current page.\loc) {\includegraphics[width=7cm]{Figures_booklet/SOM_12_Nuclear_physics_figures/SOM_Nuclear_physics_Figure_1.png}}; % 8cm max height
%\end{tikzpicture}

%\endofslide 
\endofsection
\end{frame}
%%%%%%%%%%%%%%%



